\section{Introduction}
\label{sec:intro}

On October 19, 2018 the Uber autonomous vehicle crash in Arizona highlighted a fundamental challenge in open-world perception: the difficulty of reliably detecting objects outside the training distribution. The perception system identified the pedestrian several seconds before impact but alternately classified her as a vehicle, bicycle, or unknown object, ultimately suppressing the alert. This tragedy underscores why detecting Out-of-Distribution (OoD) objects is critical for safe autonomous driving.

During training, semantic segmentation models are exposed to a closed set of In-Distribution classes; however, real-world driving environments are inherently open and unpredictable. Objects such as scattered debris, unexpected animals, or fallen cargo may appear without warning, and failing to segment them can create severe safety hazards. As a result, OoD detection remains a central problem in open-world recognition and continual learning.

Historically, semantic segmentation has been dominated by pixel-based convolutional neural networks (CNNs) such as ERFNet, which assign class labels independently to each pixel through hierarchical local feature extraction. While this design enables real-time inference, CNNs rely on limited receptive fields, essentially processing images through a sliding window of local context. Recently, mask-based architectures leveraging Vision Transformers (ViTs), such as Mask2Former, have emerged as powerful alternatives. These models fundamentally reframe segmentation as a set prediction problem: instead of classifying pixels, they employ learnable query tokens that attend globally across the entire image to generate spatial masks.

We systematically evaluate ERFNet against the Encoder-only Mask Transformer (EoMT), a highly efficient ViT powered by a DINOv2 backbone. Unlike traditional ViTs that require extensive task-specific pretraining, DINOv2 provides robust visual features learned through self-supervised learning on diverse, uncurated data, making it particularly well-suited for open-world scenarios where training data cannot cover all possible objects.

Despite the proliferation of specialized benchmarks (SMIYC, Fishyscapes, Road Anomaly), existing evaluations remain constrained to global pixel-wise metrics such as mean Intersection over Union (mIoU) and False Positive Rate at 95\% True Positive Rate (FPR95). While these aggregated statistics provide high-level overviews, they obscure critical localized failure modes: where exactly do models fail? On which object scales? At what depth bands? Near which semantic boundaries? Without this fine-grained diagnostic lens, practitioners cannot design targeted mitigation strategies, and the structural trade-offs between pixel-based CNNs and query-based ViTs remain opaque.

To extract anomaly scores without retraining, we apply a suite of post-hoc evaluation methods during inference. These methods operate directly on the model's output logits to estimate uncertainty: Maximum Softmax Probability (MSP) treats low confidence as anomalous; MaxLogit uses unnormalized activations; MaxEntropy measures predictive uncertainty; Temperature Scaling recalibrates confidence estimates. Crucially, we also evaluate Rejected-by-All (RbA), a mask-specific scoring mechanism that exploits the query-based structure of ViTs by identifying pixels that no object query claims, uniquely suited to mask architectures.

Our primary contribution is a rigorous diagnostic analysis that moves beyond global metrics to pinpoint exactly where, at what scales, and near which semantic classes critical errors occur. We reveal shared architectural weaknesses, severe Boundary Uncertainty near thin structures and depth-dependent Vertical Blindness).

Crucially, we uncover diametrically opposed scale-dependent failure modes: CNNs exhibit a Scale Paradox where performance degrades catastrophically on massive anomalies (FNR increases from 64.8\% on small objects to 90.1\% on large ones), while ViTs suffer a Scale Gap, effectively acting as low-pass filters that miss small debris (FNR 51.2\% on small objects). These discoveries expose fundamental architectural limitations that generic optimization cannot overcome.

Driven by these empirical findings, we propose Smart Data Augmentation, a targeted Cut-and-Paste Outlier Exposure that addresses each identified vulnerability through data-driven priors. Our scale prior combats both the CNN's paradox and the ViT's gap by enforcing 60\% small and 30\% medium objects during training. The categorical prior intentionally overlaps anomalies with weak ID classes (Pole, Fence) to sharpen decision boundaries.